%======================================================================
\section{Introduction}\label{sec:intro}

A formula $\alpha$ is \kl{minimal} in a \kl{logic} $L$~\cite{Kom05} if
$\alpha\in L$ and no formula \kl{strictly more general} than $\alpha$ with
respect to the \kl{substitution preorder} belongs to $L$. For example, Peirce's law
$((a\to b)\to a)\to a$ is \kl{minimal} in classical logic $\CL$, whereas
$(a\to b)\to(a\to b)$ is not, since it is a proper substitution instance of
the classical tautology $p\to p$.

Since \kl{minimality} depends on the syntactic shape of a formula,
transformations of axioms into equivalent ones need not preserve it, and it
is not clear a priori whether an axiomatization by \kl{minimal} formulas
can be brought into a normal form, for instance one of bounded \kl{order},
without losing \kl{minimality}.

This question arises in the problem of Komori and
Kashima~\cite{Kom05,Kas05}, who asked which \kl{logics} can be axiomatized by
formulas \kl{minimal} in $\CL$, and in particular whether every intermediate
logic axiomatizable in this way is either intuitionistic logic $\IL$ or
$\CL$. Nakamura and Matsuda~\cite{NM21} answered the latter question in the
negative: they exhibited a formula $G'$ of \kl{order} $5$ that is \kl{minimal} in
$\CL$ and axiomatizes Gödel--Dummett logic $\LC$ over $\IL$, and they asked
whether a counterexample of \kl{order} $4$ exists \cite[Open Problem~2]{NM21}.
\plot{Check the formulation of the problem against \cite{Kom05},
\cite{Kas05} and \cite[Sect.~1]{NM21}.}

\subsection{Main results}

We show that a translation $\St$ based on \kl{extension variables}
(Definition~\ref{def:S}) preserves \kl{minimality}. It replaces every compound
subformula occurrence $u$ of $\alpha$ by a fresh variable $\ev{u}$ and records
the meaning of $\ev{u}$ in a premise $\dprem{u}$, a one-way definition whose
direction depends on the \kl{polarity} of $u$:
\[
\St(\alpha)=\dprem{u_1}\to\dots\to \dprem{u_m}\to \ev{\rt} .
\]
The translation is a variant of the polarized translation $\varphi^+$ of
Jeřábek~\cite[Def.~3.6]{Jer16}, obtained essentially by removing the
conjunctions ``$\wedge1$'' with which Jeřábek wraps each definition. Since
each premise $\dprem{u}$ contains exactly two implications, $\St(\alpha)$ has
\kl{order} at most $4$, and it can be computed in linear time.

\begin{theorem}[Main Theorem]\label{thm:main}
Let $L$ be a \kl{logic} containing the identity axiom $\axI=a\to a$. If $\alpha$ is
\kl{minimal} in $L$ and $\St(\alpha)\in L$, then $\St(\alpha)$ is \kl{minimal} in
$L$.
\end{theorem}

The hypothesis $\St(\alpha)\in L$ is satisfied for every $\alpha\in L$, and
moreover $L\lplus \alpha=L\lplus \St(\alpha)$, whenever $L$ contains the linear
implicational logic $\BCI$ (Corollary~\ref{cor:many-logics}).
Consequently, $\St$ provides a minimality-preserving normal form for
axioms: over $\BCI$, every axiomatization by $\CL$-minimal formulas
can be replaced by one of \kl{order} at most $4$
(Corollary~\ref{cor:normal-form}). The bound $4$ is optimal, because
classical tautologies of \kl{order} at most $3$ are already intuitionistic
theorems (Proposition~\ref{prop:order3}). Theorem~\ref{thm:main}, in the
generality stated above, and Corollary~\ref{cor:many-logics} have been
formalized in Lean~4 (Appendix~\ref{sec:lean}).

The hypotheses on the base logic cannot be weakened substantially. The
axiom $\axI$ cannot be dropped: there is a \kl{logic} containing $\BC$ but
not $\axI$ in which the conclusion of Theorem~\ref{thm:main} fails
(Proposition~\ref{prop:needI}). Nor can $\BCI$ be replaced by
$\CK$ or $\BK$ in Corollary~\ref{cor:many-logics}
(Appendix~\ref{app:weak}).

\subsection{Application to the Komori--Kashima problem}

Applied to $G'$, the translation yields an explicit axiom $\St(G')$ of
$\LC$ of order $4$ that is \kl{minimal} in $\CL$ (Corollary~\ref{cor:LC}). This
answers \cite[Open Problem~2]{NM21} affirmatively. Since
Corollary~\ref{cor:normal-form} holds over $\BCK$, it also applies
to Komori's original formulation of the problem over $\BCK$
\cite[footnote~1]{NM21}. Moreover, in the study of which logics are
axiomatizable by $\CL$-minimal formulas \cite[Open Problem~1]{NM21}, one
may restrict attention to formulas of order $4$
(Corollary~\ref{cor:normal-form} and Proposition~\ref{prop:order3}).

\subsection{Substructural logics}

The translation and Theorem~\ref{thm:main} extend to languages with
fusion, lattice connectives and constants. The analogue of
Corollary~\ref{cor:many-logics} holds over $\FLe$ in the
multiplicative language and over $\FLew$, in particular over
intuitionistic logic, in the full language. Over $\FLe$ with
additive connectives, however, neither $\St$ nor Jeřábek's $\varphi^+$
serves our purpose (Section~\ref{sec:fl}).

\subsection{Overview of the proof}

The proof of Theorem~\ref{thm:main} is purely syntactic; it uses only
closure under substitution and modus ponens, together with the axiom $\axI$.
It proceeds in two steps. First, a formula of $L$ is \kl{minimal} in $L$ if and
only if no \kl{one-step generalization} of it belongs to $L$, where a
\kl{one-step generalization} is obtained by replacing some occurrences of a
single subformula by a fresh variable (Proposition~\ref{cor:criterion}).
Second, every \kl{one-step generalization} $\alpha'$ of $\St(\alpha)$ can be
decoded into a \kl{one-step generalization} $\beta$ of $\alpha$: there is a
substitution $\theta$ such that
\[
\theta\alpha'=(\gamma_1\to\gamma_1)\to\dots\to(\gamma_k\to\gamma_k)\to\beta
\]
(Lemma~\ref{lem:key}). Hence, if $\alpha'$ belonged to $L$, then so would
$\beta$, contradicting the \kl{minimality} of $\alpha$. Polarity plays no role
in this argument; it is needed only to ensure that $\St(\alpha)\in L$.
Example~\ref{ex:decode} illustrates the decoding for Peirce's law.

\subsection{Related work}\label{sec:related}

The translation $\St$ is built from well-known ingredients. Extension
variables go back to
Tseitin~\cite{Tse68}, and Statman~\cite{Sta79} used them to reduce
intuitionistic propositional logic to its implicational fragment (see also
Švejdar~\cite{Sve03}). One-way definitions chosen according to \kl{polarity}
appear in the Plaisted--Greenbaum translation~\cite{PG86} and in Mints'
normal form~\cite{Min90}, and Rybakov~\cite{Ryb97} uses \kl{extension variables}
in the study of admissible rules.

Our translation is closest to that of Jeřábek~\cite{Jer16}. His translation $\varphi^+$ introduces a fresh
variable for every occurrence of a non-variable subformula, uses the same
polarized one-way definitions, and wraps each of them as $\dprem{u}\wedge1$. The
translation $\St$ differs from $\varphi^+$ only in that it omits these
wrappers, curries the product of the definitions into an iterated
implication, and treats constants as atoms. With the wrappers,
$\FLe\lplus \varphi=\FLe\lplus \varphi^+$ holds for every formula
$\varphi$ in the full language \cite[Thm.~3.8]{Jer16}; Jeřábek uses this to
show that the substructural hierarchy of Ciabattoni, Galatos and
Terui~\cite{CGT08} collapses to the level $\mathcal N_3$ over
$\FLe$. The main point of the present paper is the observation that,
without the wrappers, the translation preserves \kl{minimality}; to our
knowledge, this has not been observed before. The wrappers destroy
\kl{minimality} (Section~\ref{sec:fl}).
\plot{Check the exact forms used in \cite{Sta79}, \cite{Sve03}, \cite{Min90}.}

\subsection{Organization of the paper}

Section~\ref{sec:prelim} fixes notation and terminology,
Section~\ref{sec:translation} defines the translation $\St$, and
Section~\ref{sec:proof} proves Theorem~\ref{thm:main}.
Section~\ref{sec:normal-form} establishes the basic properties of $\St$
over $\BCI$ and derives the normal form, and Section~\ref{sec:kk} applies
it to the Komori--Kashima problem. Section~\ref{sec:remarks} shows that the
axiom $\axI$ cannot be dropped from Theorem~\ref{thm:main}, and
Section~\ref{sec:fl} extends the results to substructural logics.
Section~\ref{sec:conclusion} concludes with open problems.
Appendix~\ref{app:criterion} proves Proposition~\ref{cor:criterion},
Appendix~\ref{sec:lean} describes the Lean formalization, and
Appendix~\ref{app:weak} shows that $\axB$ and $\axC$ cannot be dropped from
Corollary~\ref{cor:many-logics}.
